\section{结论}\label{sec:conclusion}
本文把AEIC结构化报告作为Leader--Follower实验的可观测接口，并通过行动前公开合同与真实资源路径绑定构造语义共识机制。公开Static AEIC规则显著改变主后端的资源行动，联合执行配置下的集体准确率同时改善，仅结果奖励没有产生同样结果。受控信息实验支持AEIC信号的行为测量价值，但字段消融未证明单一字段独立必要，现有质量增益也不能归因于AEIC格式本身。剂量响应、稳健性实验和冻结溯源支持这一有边界的解释；留出后端失败进一步表明有效合同必须针对解释它的模型策略校准。后续识别顺序应是先验证报告接口，再分离合同与资源效应，最后进行有限合同选择。

因此，本文完成的是固定公开合同下Follower经验响应的识别，而不是一般Stackelberg均衡证明。其可复用贡献包括统一报告Schema、先选择后执行的资源约束、过程信号结算、跨轮信誉状态以及可复算Trace。未来工作将在冻结任务上估计奖励--成本响应面，定义包含质量、资源和失败风险的Leader效用，并只在有限合同空间中选择经验最优方案；任何新合同均需在独立任务和新模型后端上重新校准与验证。
